\documentclass{article}
\usepackage{iclr2027_conference,times}
\input{math_commands.tex}
\usepackage[hidelinks]{hyperref}
\usepackage{url}
\usepackage{amsmath,amssymb,amsthm}

\newtheorem{theorem}{Theorem}
\newtheorem{lemma}{Lemma}
\newtheorem{corollary}{Corollary}
\newcommand{\Expo}{E}
\newcommand{\Budget}{B}

\title{Three Scoped Contraction Results\\
Quadratic, Convex, and Affine Scalar Regimes}
\author{Anonymous}

\begin{document}
\maketitle

\begin{abstract}
This theory-only note records three separate, scoped identities. A fixed
quadratic has an exact matrix-product solution whose first-order log
contraction depends on cumulative learning-rate exposure, with a finite-step
squared-step remainder. Under explicit pointwise ambient Hessian bounds on
every minimizer-to-iterate segment, gradient descent on a twice continuously
differentiable objective contracts full parameter error; a faster directional
rate also requires an invariant-subspace condition. An affine scalar
conditional-mean recursion has an exact product law and an exponential upper
bound. The statements do not compare two paths and are not laws for logistic
regression, shuffled SGD, CNNs, nonconvex optimization, or real label noise.
The package makes no empirical, deployment, or novelty claim.
\end{abstract}

\section{Introduction}
\label{sec:intro}

This paper is a scoped formal exposition of three product/contraction
calculations. It does not define an alternative trajectory, a reference
trajectory, or a quantity measuring their difference. Accordingly, it does
not establish a comparator-path result or data attribution. Those topics
motivate possible future research but are not conclusions of the results below.

The three results have deliberately separate domains. T1 is a fixed-quadratic
product identity. T2 is a smooth-convex full-norm contraction theorem under
pointwise ambient Hessian bounds. T3 is an affine scalar conditional-mean
identity. They share a product form but no theorem connects their state spaces
or consequences. Neither T1 nor T3 establishes a law for nonlinear SGD, and
T2 does not transfer automatically to arbitrary neural networks.

\begin{center}
\small
\begingroup
\setlength{\tabcolsep}{3pt}
\begin{tabular}{p{0.12\linewidth}|p{0.23\linewidth}|p{0.29\linewidth}|p{0.25\linewidth}}
\textbf{Result} & \textbf{State and assumptions} & \textbf{Conclusion} & \textbf{Not supplied} \\
\hline
T1 & Fixed SPD quadratic; $0\leq\eta_tL<1$ & Exact terminal matrix product and a log-product remainder & A changed objective, a comparator path, or nonlinear transfer \\
T2 & Gradient descent; pointwise $\mu I\preceq\nabla^2L_A\preceq LI$ on each minimizer--iterate segment & Full-norm contraction to $w_A^\star$; directional rate only with invariance & A comparison to another trajectory \\
T3 & Affine scalar recursion with conditional-zero-mean innovation & Exact mean product and exponential upper bound & A convex-optimization or general-SGD statement
\end{tabular}
\endgroup
\end{center}

\paragraph{Scope.}
The frozen public baseline contained historical numerical material without
co-located public code, output, configuration, seed, and preregistration
artifacts. That material is not part of this submission source. A separate,
uncompiled archive points to the immutable checkpoint containing it. No
unexecuted empirical policy is part of the active paper.

\section{Notation}
\label{sec:setup}

Write $\eta_t\geq0$ for a deterministic step size and
$E_{s:T}=\sum_{t=s}^{T-1}\eta_t$ for tail exposure. The symbol $A$ in T2
only labels the displayed objective $L_A$; it does not encode a data block or
a reference path. ``Contraction'' below always means the
particular norm or expectation displayed in its own theorem, not accuracy, a
representation, or a deployment outcome.

\section{Three bounded results}
\label{sec:theory}

\subsection{T1: fixed-quadratic exposure product}
\label{sec:t1}

\begin{lemma}[T1]
\label{thm:t1}
Let
\[
 Q(w)=\tfrac12(w-w^\star)^\top H(w-w^\star),\qquad
 0\prec\mu I\preceq H\preceq LI,
\]
and run $w_{t+1}=w_t-\eta_t\nabla Q(w_t)$ with
$0\leq\eta_tL<1$. Then
\begin{equation}
 w_T-w^\star=\prod_{t=0}^{T-1}(I-\eta_tH)(w_0-w^\star).
\label{eq:t1product}
\end{equation}
For an eigenvalue $\lambda_j$ of $H$, set
$P_j=\prod_t(1-\eta_t\lambda_j)$,
$E=\sum_t\eta_t$, and $q_j=\max_t\eta_t\lambda_j<1$. Then
\begin{equation}
 0\leq-\log P_j-\lambda_jE
 \leq\frac{\lambda_j^2\sum_t\eta_t^2}{2(1-q_j)}.
\label{eq:t1log}
\end{equation}
\end{lemma}

\begin{proof}
Since $\nabla Q(w)=H(w-w^\star)$, one step gives
$w_{t+1}-w^\star=(I-\eta_tH)(w_t-w^\star)$, and iteration proves
\eqref{eq:t1product}. Diagonalize the symmetric positive-definite $H$.
On an eigenvector with eigenvalue $\lambda_j$, the multiplier is $P_j$.
For $x\in[0,q_j]$,
\[
 0\leq-\log(1-x)-x
   =\sum_{k\geq2}\frac{x^k}{k}
 \leq\frac{x^2}{2(1-q_j)}.
\]
Apply this inequality to $x=\eta_t\lambda_j$ and sum over $t$.
This proves \eqref{eq:t1log}.
\end{proof}

T1 says that equal exposure yields the same first-order log contraction in
this fixed quadratic. It does not say that arbitrary finite schedules with
equal exposure have identical terminal iterates: the squared-step term
quantifies the permitted discrepancy. Replacing a step $\eta$ by $K$ steps of
size $\eta/K$ preserves exposure, whereas repeating $K$ original steps changes
exposure and is not a finite-budget invariance.

\subsection{T2: ambient-Hessian gradient-descent contraction}
\label{sec:t2}

\begin{theorem}[T2]
\label{thm:t2}
Let $L_A$ be twice continuously differentiable with unique minimizer
$w_A^\star$ and $\nabla L_A(w_A^\star)=0$. Suppose that for every
$t=s,\ldots,T-1$ and every $r\in[0,1]$,
\begin{equation}
 \mu I\preceq \nabla^2L_A\!\left(w_A^\star+r(w_t-w_A^\star)\right)
 \preceq LI,
 \qquad 0<\mu\leq L.
\label{eq:t2ambient}
\end{equation}
If gradient descent on $L_A$ is run from the arbitrary starting iterate $w_s$
with $0\leq\eta_t\leq1/L$, then
\begin{equation}
 \lVert w_T-w_A^\star\rVert_2
 \leq\prod_{t=s}^{T-1}(1-\eta_t\mu)
      \lVert w_s-w_A^\star\rVert_2
 \leq e^{-\mu E_{s:T}}\lVert w_s-w_A^\star\rVert_2.
\label{eq:t2full}
\end{equation}
If, in addition, $e\ne0$, $\lambda_e>0$, and $\operatorname{span}(e)$ is invariant for every segment
average Hessian in the proof and its eigenvalue on $e$ is at least
$\lambda_e$, then
\begin{equation}
 |e^\top(w_T-w_A^\star)|
 \leq e^{-\lambda_eE_{s:T}}|e^\top(w_s-w_A^\star)|.
\label{eq:t2direction}
\end{equation}
\end{theorem}

\begin{proof}
For each $t\geq s$, the fundamental theorem of calculus gives
\[
 \nabla L_A(w_t)-\nabla L_A(w_A^\star)
 =\bar H_t(w_t-w_A^\star),\qquad
 \bar H_t=\int_0^1\nabla^2L_A(w_A^\star+
        r(w_t-w_A^\star))\,dr .
\]
Integrating the pointwise ambient bounds in \eqref{eq:t2ambient} implies
$\mu I\preceq\bar H_t\preceq LI$. Therefore
\[
 w_{t+1}-w_A^\star=(I-\eta_t\bar H_t)(w_t-w_A^\star).
\]
All eigenvalues of $I-\eta_t\bar H_t$ lie in
$[0,1-\eta_t\mu]$ because $\eta_t\leq1/L$. Hence its spectral norm is at most
$1-\eta_t\mu$. Multiply these bounds from $s$ through $T-1$, then use
$1-x\leq e^{-x}$ to obtain \eqref{eq:t2full}.

For the directional statement, the added invariance premise makes $e$ an
eigenvector of every $\bar H_t$. Its scalar multiplier is at most
$1-\eta_t\lambda_e$; the same multiplication and exponential bound prove
\eqref{eq:t2direction}. A Rayleigh-quotient lower bound alone does not imply
this shared invariant-subspace premise.
\end{proof}

Thus T2 is a full-parameter contraction result under the explicit pointwise
ambient Hessian condition \eqref{eq:t2ambient}. It does not compare $w_t$ to a
second trajectory or identify why $w_s$ has its value. Its directional
refinement is conditional and must not be inferred from observations in another
model class.

\subsection{T3: affine scalar conditional-mean product}
\label{sec:t3}

\begin{theorem}[T3]
\label{thm:t3}
Let $(u_t)$ be an integrable scalar process adapted to
$(\mathcal F_t)$. Fix $\lambda>0$ and deterministic step sizes satisfying
$0<\eta_t\lambda<1$. Suppose
\begin{equation}
 u_{t+1}=(1-\eta_t\lambda)u_t+\eta_t\xi_t,\qquad
 \mathbb E[\xi_t\mid\mathcal F_t]=0.
\label{eq:t3affine}
\end{equation}
Then for any $s<T$,
\begin{equation}
 \mathbb E[u_T]=\mathbb E[u_s]\prod_{t=s}^{T-1}
 (1-\eta_t\lambda),\qquad
 |\mathbb E[u_T]|
 \leq|\mathbb E[u_s]|e^{-\lambda E_{s:T}}.
\label{eq:t3product}
\end{equation}
\end{theorem}

\begin{proof}
Taking conditional expectations in \eqref{eq:t3affine} and then total
expectations yields
\[
 \mathbb E[u_{t+1}]=(1-\eta_t\lambda)\mathbb E[u_t].
\]
Repeated substitution gives the product identity in \eqref{eq:t3product}.
Every factor is positive and $1-x\leq e^{-x}$ for $x\geq0$, so multiplying
the inequalities gives the exponential upper bound.
\end{proof}

T3 is exact for the displayed scalar recursion. It does not assert that
logistic-SGD, random reshuffling, or a nonlinear accuracy functional obeys
\eqref{eq:t3affine}.

\section{Future work}
\label{sec:empirical-status}

No empirical policy or experiment is defined or evaluated in this paper.
Possible future work could formulate an explicit comparator-path theorem or
preregister and execute a separately defined empirical study. Such work would
require its own inputs, comparator, and evidence chain; it would test practical
relevance, not serve as evidence for T1--T3.

\section{Related work}
\label{sec:related}

The locked sources on training-data/model relationships---influence functions
\citep{koh2017understanding}, datamodels \citep{ilyas2022datamodels}, TRAK
\citep{park2023trak}, and SISA \citep{bourtoule2021sisa}---are contextual
rather than antecedents for the three results. Work on SGD noise and implicit
regularization \citep{mandt2017sgd,smith2017bayesian,smith2021origin}
supports stating T3's affine conditional-mean premise instead of treating it
as a general SGD identity. The currently locked sources do not provide a
verified theorem-level antecedent comparison for all of T1--T3. We therefore
make no claim that the individual identities are novel, nor that their
juxtaposition is a nontrivial unifying result. The closest-work matrix records
this boundary and the remaining positioning gap.

\section{Limitations}
\label{sec:limitations}

The scope is mathematical rather than empirical. T1 is fixed quadratic. T2
requires the pointwise ambient spectral Hessian bounds
\eqref{eq:t2ambient} on every relevant minimizer--iterate segment, a stable
step-size bound, and a separate invariant-subspace condition for its directional
refinement. T3 requires the affine scalar recursion and conditional-zero-mean
innovation. None of these assumptions automatically holds for logistic
training, shuffled SGD, neural networks, noisy labels, or data-selection
systems. The public package also lacks a same-version experiment evidence
chain; remote candidate results are not attributed to this manuscript.

The contribution ceiling is also substantive: without a theorem that connects
the three regimes, a verified theorem-level distinction from direct antecedents,
or new empirical evidence, this package remains a scoped exposition rather
than an established research advance.

\section{Conclusion}

T1--T3 establish exact or bounded product relations only within their separate
stated regimes. They do not establish a comparator-path identity or a
deployment recommendation. A future paper may add a common formal model or
empirical evidence; this one does not.

\section*{Reproducibility Statement}

This self-contained source package includes its consumed TeX dependencies and
bibliography. The controlled build record pins the compiler and BibTeX
versions, hashes every consumed source input (including this statement), and
binds a clean rebuild to the archived PDF. Formal assumptions, statements, and
proofs are active in this source. No unexecuted protocol or remote-only result
is reported as an observation.

\section*{Ethics Statement}

This revision reports no newly run human-subject or dataset study. Any future
execution must verify dataset licenses, versions, and use conditions, and the
policy must not access clean, protected, test, or oracle labels at decision
time.

% AI Use Statement — ICLR 2027 submission.  This statement is deliberately
% revision-specific and supersedes inherited agent/workflow descriptions.

\section*{AI Use Statement}
\label{sec:ai-use}

This manuscript and its audit package were prepared with substantial assistance
from LLM-based coding agents, including Codex and earlier multi-agent research
workflows, under human supervision. The use of such systems included drafting,
revision, repository inspection, provenance and citation checks, build checks,
and internal review planning. No LLM-generated assertion was treated as
evidence without an identified source artifact or an explicit ``planned'' label.

\paragraph{Tasks performed with AI assistance.}
\begin{itemize}
  \item Drafting and revising prose, equations, figures, tables, and this
  disclosure; generating alternative paper stories and revision plans.
  \item Inspecting the public baseline and a separately frozen remote candidate;
  computing hashes and diffs; checking source-package dependencies; and running
  safe, read-only consistency/configuration verifiers.
  \item Searching for and cross-checking primary or official bibliographic
  metadata, then building a claim--evidence map and closest-work matrix.
  \item Designing (but not submitting or executing) the planned matched-path
  experiment and auditing its data, environment, launcher, comparator, seed,
  and output-schema readiness.
  \item Conducting or assisting internal manuscript audits. Such internal
  reviews are not peer review and do not establish experimental results,
  theorem correctness, or acceptance prospects.
\end{itemize}

\paragraph{Human responsibility and controls.}
\begin{itemize}
  \item Human supervisors specify the task, authorization boundaries, resource
  limits, and submission decisions. This revision did not authorize a GPU,
  remote, or formal experimental run.
  \item Human authors must verify all final claims, citations, proofs, data-use
  permissions, and venue compliance, and take full responsibility for the
  submission and any remaining errors.
\end{itemize}

The current package intentionally distinguishes the public manuscript version
from a different remote candidate. Remote scripts and results are not used as
evidence for public-version claims unless a claim-level reconciliation closes
that gap. Accordingly, this disclosure does not imply that any historical
agent-produced numerical claim has been independently reproduced.


\bibliographystyle{iclr2027_conference}
\bibliography{references}

\appendix
\section{Proof-scope checklist}

T1 uses a fixed symmetric positive-definite Hessian and its eigendecomposition.
T2 uses the pointwise ambient bounds \eqref{eq:t2ambient}, whose integration
bounds every segment-average Hessian between $\mu I$ and $LI$; its directional
line is conditional on simultaneous invariance. T3 takes conditional
expectations of the stated scalar recursion. These are the only proof
dependencies used by the active manuscript. Historical figures, tables,
scripts, and numerical claims are not imported, hidden, or compiled here.

\end{document}
